\begin{explanationbox}
	\textbf{\textcolor{CExplanation}{一句话直觉}}

	动点$P$到定点$F$的距离与到定直线$\ell$的距离恰好相等，此时$P$的轨迹就是抛物线。

	\medskip
	\textbf{\textcolor{CExplanation}{核心拆解}}
	\begin{description}[style=nextline, leftmargin=2.8em, labelsep=0.5em, itemsep=0.45em]
		\item[\textbf{要点一（定义组成）：}] 条件中涉及两个关键元素：定点$F$（焦点）和定直线$\ell$（准线），且$F$不在$\ell$上。
		\item[\textbf{要点二（等距条件）：}] 动点$P$到$F$的距离$PF$等于$P$到$\ell$的垂线段长度$PH$。这是判断轨迹的核心等式。
		\item[\textbf{要点三（轨迹结论）：}] 满足该等式的所有点$P$构成一条开口朝向$\ell$的抛物线，且$F$和$\ell$分别是其焦点和准线。
	\end{description}

	\medskip
	\textbf{\textcolor{CExplanation}{几何本质（统一圆锥曲线）}}

	抛物线是圆锥曲线中离心率$e=1$的特例。到定点与定直线的距离之比等于1，正是抛物线的本质特征。

	\medskip
	\textbf{\textcolor{CExplanation}{代数意义（坐标化条件）}}

	在恰当坐标系下，设$F\bigl(\frac{p}{2},0\bigr)$，$\ell:x=-\frac{p}{2}$，则$PF=PH$可化为$y^2=2px$，正是抛物线标准方程。

	\medskip
	\textbf{\textcolor{CExplanation}{考点价值（常考转化）}}
	\begin{itemize}[leftmargin=*, itemsep=0.5em, topsep=0.4em]
		\item \textbf{考法一（轨迹判定）：} 给出动点满足的距离关系，判定是否为抛物线。
		\item \textbf{考法二（距离转换）：} 利用定义将抛物线上的点到焦点的距离转化为到准线的距离，反之亦然，简化计算。
	\end{itemize}

	\medskip
	\textbf{\textcolor{CExplanation}{顿悟点}}

	\textbf{“把$PF$换成$PH$”或反向替换，能避免直接使用两点间距离公式带来的根号运算，是简化抛物线问题的核心技巧。}

	\medskip
	\textbf{\textcolor{CExplanation}{使用场景}}
	\begin{itemize}[leftmargin=*, itemsep=0.5em, topsep=0.4em]
		\item \textbf{场景一（轨迹判断）：} 若题目给出“$PF=d(P,\ell)$”的条件，直接判断动点轨迹是抛物线。
		\item \textbf{场景二（距离计算）：} 已知抛物线上一点坐标或焦半径，利用定义快速求点到焦点或准线的距离。
	\end{itemize}
\end{explanationbox}
